\documentclass[10pt]{article}
\usepackage[utf8]{inputenc}
\usepackage[T1]{fontenc}
\usepackage[a4paper, left=2.54cm, right=2.54cm, top=3cm, bottom=3cm]{geometry}
\usepackage[spanish]{babel}
\usepackage{amsmath, amssymb}
\usepackage{multicol}
\usepackage{enumitem}
\usepackage{fancyhdr}
\usepackage{titlesec}
\usepackage{xcolor}
\usepackage{tcolorbox}
\usepackage{booktabs}
\usepackage{array}
\usepackage{tikz}

\tcbuselibrary{skins}

% ── Colores ───────────────────────────────────────────────────────────────────
\definecolor{azulprim}{RGB}{10, 60, 130}
\definecolor{azulclaro}{RGB}{220, 232, 250}
\definecolor{verdeprim}{RGB}{0, 110, 60}
\definecolor{verdelight}{RGB}{215, 245, 225}
\definecolor{rojoprim}{RGB}{180, 30, 30}
\definecolor{naranja}{RGB}{200, 90, 0}
\definecolor{morado}{RGB}{90, 30, 140}
\definecolor{grisclaro}{RGB}{248, 248, 250}

% ── Encabezado ────────────────────────────────────────────────────────────────
\pagestyle{fancy}
\fancyhf{}
\fancyhead[L]{\small\textcolor{azulprim}{\textbf{CTA -- Clase 1} $|$ Ejercicios Propuestos}}
\fancyhead[R]{\small\textcolor{azulprim}{\today}}
\fancyfoot[C]{\small\textcolor{azulprim}{\thepage}}
\renewcommand{\headrulewidth}{1pt}
\renewcommand{\headrule}{\hbox to\headwidth{\color{azulprim}\leaders\hrule height\headrulewidth\hfill}}

% ── Bloque título de parte ────────────────────────────────────────────────────
\newcommand{\parte}[2]{%
  \begin{tcolorbox}[colback=#1, colframe=#1, arc=4pt,
    left=8pt, right=8pt, top=3pt, bottom=3pt,
    before skip=14pt, after skip=6pt]
    \textbf{\color{white}\large #2}
  \end{tcolorbox}%
}

\newcommand{\ptos}[1]{\hfill\textcolor{naranja}{\small(#1~pts.)}}

% Líneas de respuesta
\newcommand{\lineas}[1]{%
  \par\vspace{2pt}%
  \foreach \i in {1,...,#1}{%
    \noindent\rule{\linewidth}{0.3pt}\par\vspace{3pt}%
  }%
}

\setlength{\parskip}{0.4em}
\setlength{\parindent}{0em}
\setlength{\columnsep}{1.1cm}

% ─────────────────────────────────────────────────────────────────────────────
\begin{document}
% ─────────────────────────────────────────────────────────────────────────────

% ── Ficha del alumno ──────────────────────────────────────────────────────────
\begin{tcolorbox}[colback=azulclaro, colframe=azulprim, arc=4pt]
  \begin{center}
    {\large\bfseries\color{azulprim} CIENCIA, TECNOLOGÍA Y AMBIENTE -- CLASE 1}\\[3pt]
    {\normalsize\bfseries Materia, Energía, Sistema Solar y Universo}\\[3pt]
    {\small Ejercicios Propuestos}
  \end{center}
  \vspace{4pt}
  \begin{tabular*}{\linewidth}{@{}p{5.5cm}p{6.5cm}p{3cm}@{}}
    \textbf{Apellidos y nombres:} \underline{\hspace{4.5cm}} &
    \textbf{Grado/Sección:} \underline{\hspace{2cm}} &
    \textbf{Puntaje:} \underline{\hspace{1.8cm}} \\[6pt]
    \textbf{Fecha:} \underline{\hspace{4.8cm}} &
    \textbf{Docente:} \underline{\hspace{3cm}} & \\
  \end{tabular*}
\end{tcolorbox}

\vspace{0.2cm}

% ═════════════════════════════════════════════════════════════════════════════
\parte{azulprim}{PARTE I: Materia, Propiedades y Estados (30 puntos)}
% ═════════════════════════════════════════════════════════════════════════════

\begin{multicols}{2}
\begin{enumerate}[label=\textbf{\arabic*.}, leftmargin=*, itemsep=0.55cm]

  \item Define con tus propias palabras qué es la \textbf{materia} e indica sus dos características fundamentales.\ptos{2}
  \lineas{3}

  \item Clasifica las siguientes propiedades en \textbf{generales} o \textbf{específicas} de la materia, marcando con una X en el cuadro correspondiente:\ptos{3}
  
  \begin{center}
  \begin{tabular}{lcc}
  \toprule
  Propiedad & General & Específica \\
  \midrule
  Masa & \quad\fbox{\phantom{X}} & \quad\fbox{\phantom{X}} \\
  Ductilidad & \quad\fbox{\phantom{X}} & \quad\fbox{\phantom{X}} \\
  Volumen & \quad\fbox{\phantom{X}} & \quad\fbox{\phantom{X}} \\
  Magnetismo & \quad\fbox{\phantom{X}} & \quad\fbox{\phantom{X}} \\
  Inercia & \quad\fbox{\phantom{X}} & \quad\fbox{\phantom{X}} \\
  Solubilidad & \quad\fbox{\phantom{X}} & \quad\fbox{\phantom{X}} \\
  \bottomrule
  \end{tabular}
  \end{center}

  \item Un bloque de madera tiene una masa de $600\,\mathrm{g}$ y un volumen de $750\,\mathrm{cm}^3$. Calcula su \textbf{densidad} en g/cm³ y determina si flotará en agua ($\rho_{\text{agua}}=1\,\mathrm{g/cm}^3$).\ptos{3}
  \lineas{4}

  \item Una esfera de cobre tiene masa $178\,\mathrm{g}$ y densidad $8{,}9\,\mathrm{g/cm}^3$. ¿Cuál es su volumen?\ptos{2}
  \lineas{3}

  \item Completa la tabla de cambios de estado:\ptos{4}

  \begin{center}
  \begin{tabular}{lll}
  \toprule
  Cambio & De $\to$ A & Nombre \\
  \midrule
  \phantom{XX} & Sólido $\to$ Líquido & \underline{\hspace{2cm}} \\
  \phantom{XX} & Líquido $\to$ Gaseoso & \underline{\hspace{2cm}} \\
  \phantom{XX} & Gaseoso $\to$ Líquido & \underline{\hspace{2cm}} \\
  \phantom{XX} & Líquido $\to$ Sólido & \underline{\hspace{2cm}} \\
  \phantom{XX} & Sólido $\to$ Gaseoso & \underline{\hspace{2cm}} \\
  \phantom{XX} & Gaseoso $\to$ Sólido & \underline{\hspace{2cm}} \\
  \bottomrule
  \end{tabular}
  \end{center}

  \item Explica la diferencia entre un \textbf{sólido} y un \textbf{líquido} en cuanto a forma y volumen. Da un ejemplo cotidiano de cada uno.\ptos{2}
  \lineas{3}

  \item ¿Por qué el aceite flota sobre el agua? Usa el concepto de \textbf{densidad} para justificar tu respuesta.\ptos{2}
  \lineas{3}

  \item Indica qué propiedad específica de la materia se aprovecha en cada caso:\ptos{3}
  \begin{enumerate}[label=\alph*)]
    \item Fabricar cables eléctricos de cobre: \underline{\hspace{3cm}}
    \item Hacer láminas de aluminio para envolver alimentos: \underline{\hspace{3cm}}
    \item Usar goma para fabricar bandas elásticas: \underline{\hspace{3cm}}
  \end{enumerate}

  \item Un recipiente contiene $2\,\mathrm{L}$ de agua ($\rho = 1\,\mathrm{kg/L}$) y $1\,\mathrm{L}$ de aceite ($\rho = 0{,}9\,\mathrm{kg/L}$). Calcula la masa total de los líquidos.\ptos{2}
  \lineas{3}

  \item \textbf{(Desafío)} El plasma es considerado el cuarto estado de la materia. Menciona dos lugares del universo o de la vida cotidiana donde pueda encontrarse plasma.\ptos{3}
  \lineas{3}

\end{enumerate}
\end{multicols}

% ═════════════════════════════════════════════════════════════════════════════
\parte{verdeprim}{PARTE II: Energía, Fuentes y Conservación (25 puntos)}
% ═════════════════════════════════════════════════════════════════════════════

\begin{multicols}{2}
\begin{enumerate}[label=\textbf{\arabic*.}, leftmargin=*, itemsep=0.55cm, start=11]

  \item Define \textbf{energía} y enuncia la Ley de Conservación de la Energía.\ptos{2}
  \lineas{3}

  \item Clasifica las siguientes fuentes de energía en \textbf{renovables} o \textbf{no renovables}:\ptos{3}

  \begin{center}
  \begin{tabular}{lcc}
  \toprule
  Fuente & Renovable & No renovable \\
  \midrule
  Petróleo & \fbox{\phantom{X}} & \fbox{\phantom{X}} \\
  Solar & \fbox{\phantom{X}} & \fbox{\phantom{X}} \\
  Gas natural & \fbox{\phantom{X}} & \fbox{\phantom{X}} \\
  Eólica & \fbox{\phantom{X}} & \fbox{\phantom{X}} \\
  Nuclear & \fbox{\phantom{X}} & \fbox{\phantom{X}} \\
  Hidroeléctrica & \fbox{\phantom{X}} & \fbox{\phantom{X}} \\
  \bottomrule
  \end{tabular}
  \end{center}

  \item Un objeto de $2\,\mathrm{kg}$ se mueve a $10\,\mathrm{m/s}$. Calcula su \textbf{energía cinética}:
  \[
    E_k = \frac{1}{2}mv^2
  \]
  \ptos{3}
  \lineas{3}

  \item Un libro de $0{,}5\,\mathrm{kg}$ se encuentra en un estante a $1{,}5\,\mathrm{m}$ del suelo. Calcula su \textbf{energía potencial gravitacional} ($g = 10\,\mathrm{m/s^2}$):
  \[
    E_p = mgh
  \]
  \ptos{3}
  \lineas{3}

  \item Describe la transformación de energía que ocurre en cada situación (indica el tipo de energía inicial y final):\ptos{4}
  \begin{enumerate}[label=\alph*), itemsep=4pt]
    \item Una linterna encendida: \underline{\hspace{4cm}}
    \item Un ventilador funcionando: \underline{\hspace{4cm}}
    \item Una pelota cayendo: \underline{\hspace{4.5cm}}
    \item Una fogata ardiendo: \underline{\hspace{4.5cm}}
  \end{enumerate}

  \item ¿Por qué el Perú es considerado un país con alto potencial en energías renovables? Menciona al menos \textbf{dos} fuentes disponibles en el país.\ptos{3}
  \lineas{4}

  \item Explica con un ejemplo cotidiano por qué la energía no ``desaparece'' sino que se \textbf{transforma}.\ptos{3}
  \lineas{3}

  \item \textbf{(Desafío)} ¿Cuál es la diferencia entre energía \textbf{cinética} y \textbf{potencial}? Un automóvil en la cima de una pendiente y en movimiento en la base, ¿qué tipo(s) de energía tiene en cada caso?\ptos{4}
  \lineas{4}

\end{enumerate}
\end{multicols}

% ═════════════════════════════════════════════════════════════════════════════
\parte{rojoprim}{PARTE III: Sistema Solar, Tierra, Rocas y Minerales (25 puntos)}
% ═════════════════════════════════════════════════════════════════════════════

\begin{multicols}{2}
\begin{enumerate}[label=\textbf{\arabic*.}, leftmargin=*, itemsep=0.55cm, start=19]

  \item Escribe en orden los \textbf{8 planetas} del sistema solar de menor a mayor distancia al Sol:\ptos{4}
  
  \begin{enumerate}[label=\arabic*., itemsep=5pt]
    \item \underline{\hspace{3cm}} \quad 5. \underline{\hspace{3cm}}
    \item \underline{\hspace{3cm}} \quad 6. \underline{\hspace{3cm}}
    \item \underline{\hspace{3cm}} \quad 7. \underline{\hspace{3cm}}
    \item \underline{\hspace{3cm}} \quad 8. \underline{\hspace{3cm}}
  \end{enumerate}

  \item Relaciona cada planeta con su característica:\ptos{4}
  
  \begin{center}
  \begin{tabular}{lp{5cm}}
  \toprule
  Planeta & Característica \\
  \midrule
  \underline{\hspace{1.5cm}} & El más grande del sistema solar \\
  \underline{\hspace{1.5cm}} & El más caliente (no el más cercano al Sol) \\
  \underline{\hspace{1.5cm}} & Tiene anillos visibles de hielo \\
  \underline{\hspace{1.5cm}} & Rota de lado, eje inclinado $\approx98°$ \\
  \bottomrule
  \end{tabular}
  \end{center}

  \item Nombra y describe brevemente las \textbf{cuatro capas internas} de la Tierra:\ptos{4}
  \lineas{5}

  \item ¿Cuál es la diferencia entre el \textbf{movimiento de rotación} y el \textbf{movimiento de traslación} de la Tierra? ¿Qué fenómeno produce cada uno?\ptos{3}
  \lineas{4}

  \item Clasifica las siguientes rocas en ígneas, sedimentarias o metamórficas:\ptos{3}
  
  \begin{center}
  \begin{tabular}{lp{3cm}}
  \toprule
  Roca & Tipo \\
  \midrule
  Granito & \underline{\hspace{2.5cm}} \\
  Mármol & \underline{\hspace{2.5cm}} \\
  Caliza & \underline{\hspace{2.5cm}} \\
  Basalto & \underline{\hspace{2.5cm}} \\
  Pizarra & \underline{\hspace{2.5cm}} \\
  Arenisca & \underline{\hspace{2.5cm}} \\
  \bottomrule
  \end{tabular}
  \end{center}

  \item ¿Cuál es la diferencia entre una \textbf{roca} y un \textbf{mineral}? Da un ejemplo de cada uno.\ptos{3}
  \lineas{3}

  \item \textbf{(Desafío)} Explica por qué los Andes peruanos son ricos en minerales metálicos. Relaciona tu respuesta con el tipo de rocas y la formación geológica de la región.\ptos{4}
  \lineas{4}

\end{enumerate}
\end{multicols}

% ═════════════════════════════════════════════════════════════════════════════
\parte{morado}{PARTE IV: Teorías sobre el Origen y Evolución del Universo (20 puntos)}
% ═════════════════════════════════════════════════════════════════════════════

\begin{multicols}{2}
\begin{enumerate}[label=\textbf{\arabic*.}, leftmargin=*, itemsep=0.55cm, start=26]

  \item Explica brevemente en qué consiste la \textbf{Teoría del Big Bang} e indica hace cuántos años ocurrió.\ptos{3}
  \lineas{4}

  \item Menciona y explica \textbf{dos evidencias científicas} que apoyan la Teoría del Big Bang.\ptos{4}
  \lineas{4}

  \item ¿Qué establece la \textbf{Ley de Hubble}? ¿Qué conclusión nos permite sacar sobre el estado del Universo?\ptos{3}
  \lineas{3}

  \item Ordena de menor a mayor escala las siguientes estructuras del Universo:\ptos{3}
  
  \medskip
  \noindent\fbox{\strut\ Supercúmulo\ } \fbox{\strut\ Planeta\ } \fbox{\strut\ Galaxia\ } \fbox{\strut\ Sistema estelar\ } \fbox{\strut\ Cúmulo\ }
  
  \medskip
  Orden: \underline{\hspace{7cm}}

  \item Completa el cuadro comparativo de teorías cosmológicas:\ptos{4}
  
  \begin{center}
  \begin{tabular}{p{2.8cm}p{4.5cm}}
  \toprule
  Teoría & Idea principal \\
  \midrule
  Big Bang & \\[1.2cm]
  Estado Estacionario & \\[1.2cm]
  Universo Oscilante & \\[1.2cm]
  Multiverso & \\[1.2cm]
  \bottomrule
  \end{tabular}
  \end{center}

  \item \textbf{(Desafío)} La materia ordinaria representa solo el 5\% del Universo. ¿Qué compone el 95\% restante? ¿Por qué es difícil de detectar?\ptos{3}
  \lineas{4}

\end{enumerate}
\end{multicols}

\vspace{0.4cm}
\begin{tcolorbox}[colback=grisclaro, colframe=azulprim!50, arc=3pt,
  title=\textbf{Datos útiles}, coltitle=azulprim, colbacktitle=azulclaro]
\begin{multicols}{3}
$\rho = \dfrac{m}{V}$\quad (densidad)

\columnbreak

$E_k = \dfrac{1}{2}mv^2$\quad (energía cinética)

\columnbreak

$E_p = mgh$\quad ($g = 10\,\mathrm{m/s^2}$)
\end{multicols}
\vspace{2pt}
Edad del Universo: $\approx 13\,800$ millones de años \hfill
Distancia Tierra--Sol: $\approx 150 \times 10^6$ km (1 UA)
\end{tcolorbox}

\end{document}
